\begin{proposition} \label{proposition:field_of_real_numbers_constructed_as_the_quotient_of_the_ring_of_cauchy_sequences_in_the_rational_numbers_modulo_the_null_sequences_is_indeed_a_field}
    Let $\bbR = \mathcal{C}_\bbQ / \mathcal{N}_\bbQ$ be the \CrefAndHyperrefIfExist{definition:real_numbers_as_cauchy_completion_of_rational_numbers}{real numbers} constructed as the quotient of the \CrefAndHyperrefIfExist{definition:cauchy_sequence_in_rational_numbers}{ring of Cauchy sequences in $\bbQ$} by the \CrefAndHyperrefIfExist{definition:null_sequence_in_rational_numbers}{null sequences}\CrefIfExists{proposition:the_set_of_null_sequences_of_rational_numbers_form_an_ideal_in_the_ring_of_cauchy_sequences_of_rational_numbers}. Then:
    \begin{enumerate}
        \item $\bbR$ is a \CrefAndHyperrefIfExist{definition:commutative_ring}{commutative ring with unity} under the induced addition and multiplication:
        $$[a_n] + [b_n] = [a_n + b_n], \qquad [a_n] \cdot [b_n] = [a_n b_n].$$
        \item The zero element is $[0]$ (the class of the constant sequence $(0)$) and the unity is $[1]$ (the class of $(1)$).
        \item Every nonzero element of $\bbR$ has a multiplicative inverse. In particular, $\bbR$ is a \CrefAndHyperrefIfExist{definition:field}{field}.
    \end{enumerate}
\end{proposition}

\begin{proof}
    \TODO{AI generated, verify}
    The quotient ring $\mathcal{C}_\bbQ / \mathcal{N}_\bbQ$ is a commutative ring with unity by \Cref{definition:quotient_ring_of_a_ring_by_a_two_sided_ideal}, because $\mathcal{C}_\bbQ$ is a commutative ring with unity (\Cref{proposition:the_set_of_cauchy_sequences_of_rational_numbers_form_commutative_ring}) and $\mathcal{N}_\bbQ$ is an ideal (\Cref{proposition:the_set_of_null_sequences_of_rational_numbers_form_an_ideal_in_the_ring_of_cauchy_sequences_of_rational_numbers}). The induced operations are precisely
    $$[a_n] + [b_n] = [a_n + b_n], \qquad [a_n] \cdot [b_n] = [a_n b_n],$$
    and the zero and unity are $[0]$ and $[1]$, respectively. This establishes parts (1) and (2).

    For part (3), let $[a_n] \in \bbR$ be nonzero, i.e.\ $[a_n] \neq [0]$. This means $(a_n) \notin \mathcal{N}_\bbQ$, so $(a_n) \in \mathcal{C}_\bbQ \setminus \mathcal{N}_\bbQ$. By the proof of \Cref{proposition:the_set_of_null_sequences_of_rational_numbers_form_an_ideal_in_the_ring_of_cauchy_sequences_of_rational_numbers}, there exists a Cauchy sequence $(b_n) \in \mathcal{C}_\bbQ$ such that $(a_n)(b_n) \equiv (1) \pmod{\mathcal{N}_\bbQ}$, i.e.\ $[a_n] \cdot [b_n] = [1]$. Hence $[b_n]$ is the multiplicative inverse of $[a_n]$.

    Thus every nonzero element of $\bbR$ has a multiplicative inverse, so $\bbR$ is a \Cref{definition:field}{field}.
\end{proof}